\cotefile{136}
\scanpage{204}
Endomorphismes de l'anneau différentiel $\frac12$ simplicial $A_{\ast}$

(NB. on ne demande plus de respecter la structure de $S$-algèbre, $S=A_0^{\deleted{\unsure{*}}}$, i.e. on ne demande plus que $\varphi_0$ induise l'identité sur $A_0$).

Rappelons-nous qu'un complexe de De Rham peut se construire chaque fois qu'on a un anneau comm. $S$ avec idéal $S^+$ à puiss. divisées, et
$t\in S^+$, via une \inserted{$S$-algèbre} \deleted{algèbre} différentielle \inserted{$\wedge$-graduée à p.d.} $\frac12$ simpliciale $A_\ast(S,S^+,t)$, où
\[
A_n(S,S^+,t)=S\{X_0,\ldots,X_n\}[dX_0,\ldots,dX_n]\big/\bigl(\textstyle\sum X_i-t,\sum dX_i\bigr)_{\mathrm{pd}}
\]
Si on a un hom. d'anneaux à p.d., $S\xrightarrow{u} S'$ ($u(S^+)\subset S'{}^+\ldots$) et $\lambda'\in S'$ tel que
\deleted{$\unsure{\lambda t\in S^+|t\in S^+|\lambda\in S^+}$} $u(t)=\lambda t'$, on \deleted{en} déduit un hom. \inserted{de $S$-alg. $\wedge$-graduées différentielles à p.d.}
$A_\ast(u,\lambda): A_\ast(S,S^+,t)\longrightarrow \unsure{A'_\ast|A_\ast}(S',S'{}^+,t')$, qui en degré $n$ est l'hom.
\diagram{\begin{tikzcd}[ampersand replacement=\&,row sep=small]
S\{X_0,\ldots,X_n\}[dX_0,\ldots,dX_n]/\text{---}\arrow[d]\\
S'\{X_0,\ldots,X_n\}[dX_0,\ldots,dX_n]/\text{---}
\end{tikzcd}}
donné par \deleted{$\unsure{X_0\longleftrightarrow}$} $u:\deleted{\unsure{u}}S\longrightarrow S'$ sur les scalaires,
\[
\begin{aligned}
X_0&\longmapsto\lambda X_0\quad\bigl(\text{donc }\textstyle\sum X_0-t\longmapsto\sum\lambda X_0-u(t)=\sum\lambda X_{0}-\lambda t'\\[-2pt]
&\hspace{21em}=\lambda(\textstyle\sum X_0-t')\bigr)\\
dX_0&\longmapsto\lambda dX_0.
\end{aligned}
\]
Ainsi $A_\ast(S,S^+,t)$ devient un foncteur covariant \deleted{de} la catégorie des triples $(S,S^+,t)$, où les hom sont les couples $(u,\lambda)$ \inserted{comme ci-dessus} (composés par $(u',\lambda')(u,\lambda)=(u'u,u'(\lambda)\lambda')$).

\uline{NB} En plus de sa graduation, $A_n(S,S^+,t)$ admet une filtration \inserted{croissante} naturelle, où $\mathrm{Fil}^{\unsure{n|r}}A_n(S,S^+,t)$ est \unsure{le|les} sous
module des expressions $\sum F_{i_0\ldots i_p}\,dX_{i_0}\wedge\ldots\wedge dX_{i_p}$, avec $F_{i_0\ldots i_p}\in S\{X_0,\ldots,X_n\}$ \unsure{étant|elt|élément} de ``degré'' $\leq n-\unsure{p|r}$, ce sont des sous-$S$ modules (mais pas des idéaux), on
a $\mathrm{Fil}^n\cdot\mathrm{Fil}^m\subset\mathrm{Fil}^{n+m}$, $\mathrm{Fil}^0=S$ (scalaires), $\mathrm{Fil}^n=0$ si $n<0$. Cette filtration est respectée par les $A_\ast(u,\lambda)$.
% archive sheet number: 203

\scanpage{205}\mspage{- 67 -}
\hfill IX

En effet, la dernière hypothèse faite nous ramène au cas où
$G'=A$, $u=\text{identité}$, et où $v_{\handwritten{\eta}}:G_{\handwritten{\eta}}\handwritten{\longrightarrow}A_{\handwritten{\eta}}$ est donc une projection, qu'on se propose de prolonger en une projection $v:G\handwritten{\longrightarrow}A$. Or $G$ est \uline{lisse} comme extension de deux schémas en groupes lisses, de sorte que la conclusion résulte de la définition de la propriété néronienne de $A$.

\uline{Corollaire} 5.10. \uline{Soient} $A_K$ \uline{un schéma abélien sur le corps des fonctions} $K$ \uline{d'un trait} $S$, $\handwritten{\ell}$ \uline{un nombre premier}. \uline{Lorsque} $\handwritten{\ell}=p$, \uline{on suppose} $\mathrm{car.}K=0$,
\uline{pour pouvoir disposer du théorème de TATE déjà invoqué dans} 5.8. \uline{Pour que} $A_K$ \uline{ait bonne réduction sur} $S$, \uline{il faut et il suffit que le pro}-$\handwritten{\ell}$-\uline{groupe de Barsotti-Tate} $T_{\handwritten{\ell}}(A_K)$ ``\uline{ait bonne réduction sur} $S$'', \uline{i}.\uline{e}. \uline{se prolong}e \uline{en un groupe de Ba\deleted{\unsure{t|r}}\unsure{r|t}sotti-Tate sur} $S$.

(C'est le résultat conjecturé dans [8, 4.9].)

La nécessité étant claire, il reste à prouver la suffisance ; notons d'ailleurs que le cas $\handwritten{\ell}\ne p$ a déjà été traité (2.2.9). On peut supposer $S$ hensélien (2.2.9). Supposons d'abord que $A_K$ ait réduction semi-stable sur $S$. Alors 5.8 appliqué à $F=X$ nous montre qu'on a
$T_{\handwritten{\ell}}(A_K)=T_{\handwritten{\ell}}(A_K)^{\mathrm f}$, ce qui signifie aussi que le tore maximal de $A_{\unsure{o|0}}$ est
nul, i.e. que $A_o^o$ est un schéma abélien, de sorte que la conclusion résulte de (2.2.9 (iii$^{\mathrm o}$)). Traitons maintenant le cas général. En vertu de 2.2.9 on peut supposer évidemment $S$ strictement local.
Soir $K'$ une extension galoisienne finie telle que $A_{K'}$ ait réduction semi-stable sur le normalisé $S'$ de $S$ dans $K'$ (3.6). Si $A^*$ désigne le modèle de Néron correspondant, $G=\mathrm{Gal}(K'/K)$ opère sur la situation $(A^*,S')$, donc sur
% archive sheet number: 204
% pencil number at lower centre: 385
% stamp: numbered label 385

\scanpage{206}
Considérons le cas où $t'\in S'{}^+$ est un élément régulier de $S'$ ; alors $\lambda$ tel que $u(t)=\lambda t'$ est uniquement déterminé par $u$ ; donc si on se borne aux
triples $(S,S^+\ \text{à p.d.},t)$ tels que $t$ soit régulier, et aux hom $u:(S,S^+\deleted{,t})\longrightarrow(S',S'{}^+\deleted{,t'})$ tels que $u(t)\in S.t'$, alors $A_\ast(S,S^+,t)$ devient un foncteur sur cette catégorie, à valeurs ds la cat. des
\deleted{\unsure{alg}} anneaux gradués \inserted{(alternés)} différentiels \deleted{\unsure{act|et}} filtrés.

La condition envisagée sur $(S,S^+_{\mathrm{pd}},t)$ est vérifiée en particulier si on fait $S=k\{T\}$, $S^+=(T)_{\mathrm{pd}}$, $t=T$, lorsque $k$ est sans torsion.

Revenons alors à $\displaystyle A_\ast(k)\overset{\mathrm{dfn}}{=}A_\ast(\underset{\substack{\Vert\\k\{T\}}}{S_k},S_k^+,T)$, et à la détermination des endom de $A_\ast(k)$, comme \deleted{anneau} \inserted{$k$-algèbre} différentielle $\frac12$ simpliciale (où on ne considère \deleted{\unsure{donc}} la structure de $S$-algèbre, ni la graduation, ni les p.d.\deleted{\unsure{?}}, ni de filtration a priori\ldots)
\deleted{\unsure{Bien entendu|Bien entendu que|lorsqu'on} \unsure{que si}} Soit $\unsure{G|\mathscr G}_k$ le \deleted{groupe} \inserted{monoïde} des couples $(u,\lambda)$, où $u$ est un endom. de $S$ comme $k$-algèbre à p. div., et $\lambda\in S^{\deleted{\unsure{\times|+}}}$ tel que $u(T)=\lambda\deleted{\{T\}}T$.
Comme $u\deleted{\unsure{(T)}}$ est connu par la connaissance de $u(T)$,
$G_k$ est donc $\simeq$ le \deleted{groupe} \inserted{monoïde}
\deleted{\unsure{des} $S^\times$} \inserted{d'ens. sous-jacent $S$, avec la multiplication $(\lambda'\ast\lambda)\{T\}\deleted{\unsure{T}}=\lambda'\{T\}\deleted{\unsure{T}}\{T\lambda'\{\unsure{\lambda|T}\}\}$}. On désigne par $(u_\lambda)_\ast$ l'endomorphisme de $A_\ast(k)$ défini par $\lambda\in S$ (d'où $u_\lambda u_{\lambda'}=u_{\lambda\lambda'}$)

Il est clair que $\lambda\longmapsto u_{\lambda\ast}$ est un hom de monoïdes \uline{injectif} (en fait $\lambda\longmapsto(u_\lambda)_{\unsure{1|k|\ast}}$ l'est déjà, car $(u_\lambda)_\ast(X_0)=\lambda X_0$ et $A_1(k)\simeq k\{T\}\{X_0\}$. Ceci dit,
Si $\unsure{H|\mathscr H}_k$ est le monoïde des $G_k$-endomorphismes de la $k$-algèbre différentielle $\frac12$ simpliciale $A_\ast(k)$, \deleted{est-il vrai que} $\Omega_k$ le monoïde de \uline{tous} les endomorphismes de la $k$-alg. diff. $\frac12$ simpliciale, on a
% archive sheet number: 205

\scanpage{207}\mspage{- 52 -}
\hfill IX

Comme le foncteur $X\handwritten{\longmapsto}\hat X$ (complétion formelle le long de la fibre spéciale) induit une équivalence entre la catégorie des schémas finis sur $S$, et celle des schémas formels finis sur $S$ (EGA III 4.8), nous allons identifier simplement un schéma fini sur $S$ avec le schéma formel correspondant. Avec cette identification, la définition des ``parties fixes'' nous donne aussitôt pour tout $\handwritten{\nu}\geq 0$ :
\[
({}_{\handwritten{\ell^{\nu}}}A^{\unsure{\mathrm{o}|0}})^{\mathrm f}={}_{\handwritten{\ell^{\nu}}}(\hat A^{\mathrm o})\ ,\tag*{(5.1.4)}
\]
ce qu'on peut écrire aussi comme une identité de systèmes projectifs
\[
T_{\handwritten{\ell}}(A^{\mathrm o})^{\mathrm f}=T_{\handwritten{\ell}}(\hat A^{\mathrm o})\ .\tag*{(5.1.5)}
\]
On obtient donc un sous-pro-groupe de Barsotti-Tate de la partie fixe de $T_{\handwritten{\ell}}(A^{\mathrm o})$ en prenant $T_{\handwritten{\ell}}(\hat T)\handwritten{\subset}T_{\handwritten{\ell}}(\hat A^{\mathrm o})$. On définira alors la \uline{partie torique} de $T_{\handwritten{\ell}}(A^{\mathrm o})$ comme
\[
T_{\handwritten{\ell}}(A^{\mathrm o})^{\mathrm t}\overset{\mathrm{dfn}}{=}T_{\handwritten{\ell}}(\hat T)\handwritten{\subset}T_{\handwritten{\ell}}(\hat A^{\mathrm o})=T_{\handwritten{\ell}}(A^{\mathrm o})^{\mathrm f}\ .\tag*{(5.1.6)}
\]
Donc par \unsure{.définition|définition}, on aura
\[
T_{\handwritten{\ell}}(A^{\mathrm o})^{\mathrm t}=({}_{\handwritten{\ell^{\nu}}}\hat T)_{\handwritten{\nu}\geq 0}\ ,\tag*{(5.1.7)}
\]
le $S$-groupe fini et plat $\,{}_{\handwritten{\ell^{\nu}}}\hat T$ étant lui-même décrit comme la limite inductive des $S_n$-groupes finis et plats
\[
{}_{\handwritten{\ell^{\nu}}}(T_n)\ ,\quad n\ \text{variable}\ ,\tag*{(5.1.8)}
\]
i.e. son algèbre affine est décrite comme la limite projective des algèbres affines des groupes (5.1.8). Bien entendu, on aura
\[
(T_{\handwritten{\ell}}(A^{\mathrm o})^{\mathrm t})_n=T_{\handwritten{\ell}}(T_n)\ ,\tag*{(5.1.9)}
\]
% archive sheet number: 206
% pencil number at lower centre: 370
% stamp: numbered label 370

\scanpage{208}
\[
\unsure{G|\mathscr G}_k,\ \unsure{H|\mathscr H}_k\subset\Omega_k
\]
\[
G_k\cap H_k=\{u_\lambda\mid\lambda\in S_k,\ \lambda T=T\ \text{i.e.}\ (\lambda\{T\}-1)T=0\}
\]
Si on pose $\displaystyle\lambda\{T\}=\sum_{i\geq0}\lambda_iT^{(i)}$, alors la condition sur $\lambda$ envisagée est
\[
\boxed{\lambda_0=1}\quad\deleted{\unsure{\lambda_0=1},}\qquad\boxed{(i+1)\lambda_i=0\qquad\forall i\geq1}
\]
\deleted{ce qui implique en particulier $\lambda_0=1$.}

\deleted{\unsure{En car.} 0, cela implique \unsure{plus que}} Si $k$ sans torsion, cela implique $\lambda_i=0$ pour $i\geq1$ i.e. $\lambda=1$, i.e. $u_\lambda=\mathrm{id}$, d'où $G_k\cap H_k$\inserted{$=(\mathrm{id})$} dans ce cas\unsure{,|, si} --- et dans ce cas seulement. \deleted{\unsure{\}}}
[Dans ce cas \inserted{$k$ de car. 0 i.e. une $\mathbb Q$-alg.,}, $G_k$ est aussi le groupe \inserted{$G'_k$} des tous les endom. de $S_k=k\{T\}$ comme $k$-alg. augmentée \deleted{\unsure{p.d.}} et par suite \inserted{l'hom.} \deleted{\unsure{on a alors} $\Omega$}
\diagram{\begin{tikzcd}[ampersand replacement=\&,column sep=small]
(\unsure{1|0}\arrow[r]\&H'_k\arrow[r]\&{})\ \Omega_k\arrow[r,dash]\&G'_k=\operatorname{End}_{\substack{k\text{-alg}\\\deleted{\text{augm}}}}(\underset{\substack{\Vert\\S_k}}{\underbrace{k\{T\}}})\\
\&\&{}\&G_k\arrow[u,"(\unsure{??|\simeq|22})"']\arrow[ul,bend left=48]
\end{tikzcd}}
T
donne $\Omega_k$ comme produit $\frac12$ direct de $G_k$\inserted{$\simeq G'_k$} et de $H'_k$\ldots]

\deleted{Dans le cas général, on peut dire que}
\diagram{\begin{tikzcd}[ampersand replacement=\&,column sep=small]
1\arrow[r,dash]\&H_k\arrow[r,dash]\&\deleted{\unsure{G_k}}\Omega_k\arrow[r]\&\deleted{\unsure{\Omega'_k}}G'_k=\operatorname{End}_{k\text{-alg augm.}}(k\{T\})\\
\&\&{}\&G_k\arrow[u]\arrow[ul,bend left=48,dash,"{\unsure{\phantom{\rightarrow}|\rightarrow}}"']\\
\&\&{}\&G_k\cap H_k\simeq\{u_\lambda\mid\lambda=1+\sum_{i\geq1}\lambda_i T^{(i)},\ (i+1)\lambda_i=0\}\arrow[u]\\
\&\&{}\&{}\arrow[u]
\end{tikzcd}}
% archive sheet number: 207

\scanpage{209}\mspage{- 51 -}
\hfill IX

5. \uline{Le théorème d'orthogonalité dans le cas} $\handwritten{\ell}=p$ (\uline{cas semi-stable}). \uline{Dualité des schémas abéliens} $B_{\unsure{\mathrm o|0}}$ \uline{et} $B'_{\mathrm o}$. \uline{Caractérisation de la partie fixe}. \uline{Critères de bonne réduction}, \uline{de bonne réduction essentielle et semi-stable}.

5.1. Plaçons\unsure{-|\,}nous dans le cas 2.2.3 où $S$ est hensélien et $A^{\mathrm o}$ est $\handwritten{\ell}$-divisible, de sorte que la \uline{partie fixe}
\[
T_{\handwritten{\ell}}(A^{\mathrm o})^{\mathrm f}\handwritten{\lhook\joinrel\longrightarrow}T_{\handwritten{\ell}}(A)
\]
est définie comme le pro-$\handwritten{\ell}$-groupe de Barsotti-Tate formé des $({}_{\handwritten{\ell^\nu}}A^{\mathrm o})^{\mathrm f}$. Nous allons maintenant donner une définition de la partie torique (2.3.2) de $T_{\handwritten{\ell}}(A^{\mathrm o})$, indépendante de l'hypothèse $\handwritten{\ell}\ne p$ faite dans 2.3, mais supposant en revanche $S$ \uline{complet} (*). Pour ceci, introduisons pour tout entier $n\geq0$ le schéma artinien local
\[
S_n=\operatorname{Spec}(V/\underline{\mathrm m}^{n+1})\ ,\tag*{(5.1.1)}
\]
$n$.ème voisinage infinitésimal de $s$ dans $S$, et pour tout $S$-schéma $X$ soit $X_n=X\times_S S_n$. Comme $(A^{\mathrm o})_n=(A_n)^{\mathrm o}$ est lisse sur $S_n$, il est connu (SGA 3 IX 3.6 bis) que $T_{\mathrm o}$ est la fibre spéciale d'un unique sous-tore
\[
T_n\handwritten{\subset}A_n^{\mathrm o}\ .\tag*{(5.1.2)}
\]
D'après l'unicité des $T_n$, ceux-ci se déduisent les uns des autres par les changements de base $S_{n'}\handwritten{\longrightarrow} S_n$. En d'autres termes, il définissent un \uline{sous-tore formel} $\hat T$ du groupe $\hat A^{\mathrm o}$ (dans la catégorie des schémas formels sur $S$) complété formel de $A^{\mathrm o}$ le long de la fibre spéciale :
\[
\hat T\handwritten{\subset}\hat A^{\mathrm o}\ .\tag*{(5.1.3)}
\]
\noindent\rule{6em}{.4pt}\par
(*) Cette hypothèse n'est d'ailleurs pas essentielle, comme on montre dans 6.1 ci-dessous.
% archive sheet number: 208
% pencil number at foot of upright typescript: 369
% stamp: numbered label 369

\scanpage{210}
essayons de déterminer directement les $\varphi_\ast$, endom de la $k$-alg. diff. $\frac12$ simpliciale $A_\ast(k)$.

1) Premier invariant $\varphi_0\in\operatorname{End}_{k\text{-alg}}(S_k)$
\margin{$\begin{gathered}\checkmark\\ -\\ \cdot\\ -\\ \cdot\end{gathered}$}
\deleted{$\varphi_0$} $\varphi_0$ est donné par les
\[
\varphi_0(T^i)=\text{\textcircled{$\tau_i$}}\in S_k\qquad i\geq1
\]
satisfaisant les conditions
\[
\boxed{\tau_i\tau_j=c_{ij}\tau_{i+j}}
\]
Chaque $\varphi_n$ est $\varphi_0$-semi-linéaire.

2) Deuxième invariant $\varphi_1$, endom $\varphi_0$-semi-linéaire de $S_k\{X_0\}[dX_0]$, considéré comme $S_k$-alg. différentielle. \inserted{Il doit être compatible avec les augm. $\alpha$ et $\beta$}
Plus généralement une prop. autonome :

\uline{Prop} Soit $S$, $S'$ deux anneaux comm., $\varphi_0:S\longrightarrow S'$ un hom, on s'intéresse à la détermination des hom d'anneaux différentiels
\[
\varphi_1:\underbrace{S\{X_0\}[dX_0]}_{A_S}\longrightarrow\underbrace{S'\{X_0\}[dX_{\unsure{0|n}}]}_{A_{S'}}
\]
induisant $\varphi_0$ sur $S$, et compatible avec augmentation ($X_0\longmapsto0$, $dX_0\longmapsto0$). Ils sont donnés \unsure{exactmt|exactement} par des systèmes d'éléments
\[
\begin{aligned}
\underset{\substack{\Vert\\A_i+B_i\,dX_0}}{\xi_i}&\in A_{S'}^+=\operatorname{Ker}(\alpha_{S'}:A_{S'}\longrightarrow S')\\
&A_i\in S'\{X_0\}^+,\quad B_i\in S\{X_0\}
\end{aligned}
\]
satisfaisant à $\qquad A_i=A_1^{(i)}$
$\varphi_1$ étant déterminé par les $\xi_i$ via :
\[
\begin{aligned}
\varphi_1(X_0^{(i)})&=\xi_i\\
\varphi_1(dX_0)&=d\xi_1
\end{aligned}
\]
$\varphi$ est compatible avec les graduations (resp. les p.d.) ssi $B_i=0\quad\forall i\geq1$ (resp. $B_i=A_1^{(i-1)}B_1\quad\forall i\geq2$).
% archive sheet number: 209

\scanpage{211}\mspage{- 50 -}
\hfill IX

\margin{$\begin{gathered}\cdot\\\cdot\end{gathered}$}apparaît comme un quotient de $\mathbb Q[T]/(T^{p^N}-1)$, qui est un composé direct
de corps extensions finies de $\mathbb Q$. Il en est donc de même de $E_{\mathbb Q}$. Considérant
les images de $E,u,v$ dans les corps composants $K_i$ de $E_{\mathbb Q}$, on est donc
ramené à voir que si $K=\mathbb Q(u)$ est un corps engendré par une racine
primitive $p^n$-ème de l'unité $u$, et si $n$ est un entier $\geq2$ tel que $u-1\equiv0\bmod n$ (dans l'anneau des entiers de $K$), alors on a $u=1$, ou $u=-1$ et $n=2$.
Or ceci résulte aussitôt du fait bien connu que si $u\ne1$ i.e. $m\geq1$, l'idéal $(u-1)$ est premier et l'idéal $(p)$ en est une puissance, dont l'exposant est $>1$ sauf si $p=2,m=1$ i.e. si $u=-1$.
Du fait que $(u-1)$ est premier et divise $p$ résulte en effet que $n\geq2$ ne peut diviser $u-1$ que si $n=p$, et $p$ ne divise $u-1$ que si $(u-1)=(p)$ i.e. $u=-1$, donc $n=2$.

4.7.4. Démontrons 4.7.2. Il suffit de prouver que $\Gamma$ est commutatif.
Or si $x,y\in\Gamma$, on aura $(xy)(xy)=1$, ce qui peut s'écrire aussi, puisque $x^2=1$, $y^2=1$, $xy=(xy)^{-1}=y^{-1}x^{-1}=yx$, cqfd.

\uline{Remarque} 4.7.5. La démonstration donnée 4.7.3, jointe à 4.7.2, nous
montrent que si $G$ est un Groupe commutatif divisible d'un Topos, $n\geq2$ un entier, et $\Gamma$ un groupe fini, opérant fidèlement sur $G$ en induisant
l'identité sur $\,{}_nG$, alors $\Gamma$ est réduit à l'identité si $n\geq3$, et est isomorphe à un groupe $(\mathbb Z/2\mathbb Z)^r$ si $n=2$.
% archive sheet number: 210
% pencil number at lower centre: 368
% stamp: numbered label 368

\scanpage{212}
\[
\begin{aligned}
\varphi_1(X_0^{\deleted{\unsure{i}}(i)})&=\text{\tikz[baseline=(a.base)]\node[draw,ellipse,inner sep=2pt](a){$\xi_i,\ i\geq1$};}\qquad
\boxed{\xi_i\xi_j=c_{ij}\xi_{i+j}}\quad(i,j\geq1)\\
\varphi_1(dX_0)&=d\xi_1\\
&\boxed{d\xi_i=\xi_{i-1}\,d\xi_1}\qquad\boxed{\xi_i\in\deleted{\unsure{A}}S'\{X_0\}^+}
\end{aligned}
\]
\[
\begin{aligned}
\xi_i&=A_i+B_i\,dX_0& A_i,B_i&\in S'\{X_0\}\\
dA_i&=A_{i-1}A'_1& A_i&\in S'\{X_0\}^+\\
A_i&=\xi_1^{(i)}.
\end{aligned}
\]

\uline{Remarque} On voit donc que tt $\varphi_1$ se met de façon \uline{unique} sous la forme $\psi_1\circ\bar\varphi_0$, où $\bar\varphi_0$ est l'hom can. prolongeant
$\varphi_0$ ($\bar\varphi_0(X_0^{(i)})=X_0^{(i)}$, $\bar\varphi_0(dX_0)=\deleted{\unsure{\varphi_0(}}dX_0$) et $\psi_1$ est un $S'$-endom de la $S'$-alg. diff. $A_{S'}$, compatible avec l'augmentation. ---

Il faut maintenant exprimer \inserted{(si $S=k\{T\}$, $S'=k'\{T\}$)} la compatibilité de $\varphi_1$ avec augmentation $\beta$ ($X_0\longmapsto T$, \inserted{d'où $X_0^{(i)}\longmapsto T^{(i)}$,} $dX_0\longmapsto0$), on trouve que l'on doit avoir
$\beta_{S'}(\xi_i)=\bar\varphi_0\unsure{\vert|T}(T^{(i)})$ i.e. \unsure{doit|d'où} $\beta_{S'}(A_i)=\bar\varphi_0(T^{(i)})$
ce qui \unsure{revient|se réduit|évidemment}, puisque $A_i=A_1^{(i)}$, :
\[
\left\{\begin{aligned}
\beta_{S'}(A_1)&=\bar\varphi_0(T)\\
\varphi_0(T^{(i)})&=\varphi_0(T)^{(i)}
\end{aligned}\right.
\]
ou encore :
\[
A_1\deleted{\unsure{\{X_0\}|(X_0)}}=X_0+f\deleted{\unsure{\{X_0\}|(X_0)}}
\]
où $\displaystyle f\in A_{\deleted{\unsure{k}}1}=\ k\{X_0,X_1\}$
est donné par
\[
f=\sum_{i\geq1,\,j\geq1}X_0^{(i)}X_1^{(j)}
\]
La \uline{remarque} ci-dessus se prolonge ---
% archive sheet number: 211

\scanpage{213}\mspage{- 49 -}
\hfill IX

On peut évidemment supposer $S$ strictement local. Avec les notations de 4.1 et de 4.2, prenons $K'$ le plus petit possible, i.e.
supposons que $\Gamma$ opère \uline{fidèlement} sur ${A^{*}_{\unsure{\mathrm o|0}}}^{\mathrm o}$. Par hypothèse, $\Gamma$ opère
trivialement sur $\,{}_nA(\bar K)$, et à fortiori sur $\,{}_nA(\bar K)^{\mathrm{ef}}\simeq{}_nA^*_{\mathrm o}(k)\supset{}_n{A_{\mathrm o}^{*}}^{\mathrm o}(k)$. On conclut alors par les deux lemmes suivants :

\uline{Lemme} 4.7.1. (J.P. SERRE). \uline{Soient} $G$ \uline{un schéma en groupes commutatif sur un corps} $k$, \uline{extension d'un schéma abélien par un tore}, $n$ \uline{un nombre entier} $\geq2$. \uline{Alors tout automorphisme d'ordre fini de} $G$ \uline{induisant l'identité sur} $\,{}_nG$ \uline{est réduit à l'identité si} $n\ne2$, \uline{et est l'identité ou d'ordre deux si} $n=2$.

\uline{Lemme} 4.7.2. \uline{Un groupe fini dont tout élément} $\ne1$ \uline{est d'ordre} $2$ \uline{est isomorphe à} $(\mathbb Z/2\mathbb Z)^r$, \uline{pour un} $r\geq0$ \uline{convenable}.

4.7.3. Rappelons la démonstration de 4.7.1, qui est bien connue. On note que grâce au fait que $n\,\mathrm{id}_G$ est un épimorphisme,
l'hypothèse que $u$ induise l'identité sur $\,{}_nG$ i.e. que $1-u$ induise $0$ sur $\,{}_nG$, s'exprime par la condition que l'on ait $1-u=nv$ dans l'anneau des endomorphismes de $G$.
Pour prouver la conclusion annoncée, on est ramené (en décomposant le groupe cyclique fini engendré par $u$ en ses composantes $p$-primaires, pour les divers nombres premiers $p$ divisant l'ordre de $u$)
au cas où $u$ est d'ordre $p^N$, puissance d'un nombre premier.
Soit $E$ le sous-anneau de $\operatorname{End}(G)$
engendré par $v$, qui est un quotient de $\mathbb Z[T]/(T^{p^N}-1)$, donc de type fini sur $\mathbb Z$, et qui de plus est sans torsion (ceci est conséquence formelle du fait que $G$ est divisible).
Il se plonge donc dans $E_{\mathbb Q}=E\otimes_{\mathbb Z}\mathbb Q$, qui
% archive sheet number: 212
% pencil number at lower centre: 367
% stamp: numbered label 367

\scanpage{215}
\margin{p.d. dans les anneaux \unsure{alternés}}
% margin note written sideways
\margin{Feuille paumée\ldots}
\noindent$\left[\vphantom{\rule{0pt}{6\baselineskip}}\right.$\begin{minipage}{.90\linewidth}
\uline{Théorème} Soient $S$ un anneau muni d'un idéal $I$ à puiss. divisées, $M^\ast=M^0+M^1$ un $S$-module gradué mod $2$,
$A(M^{\deleted{\unsure{0}}\ast})=\Gamma^\ast(M^0)\otimes_S\bigwedge^\ast M^1$ \inserted{$S$-algèbre} gradué par la parité du degré total, [et augmentée vers $S$ de la façon évidente $\varepsilon:A\longrightarrow S$]
$J=$ l'idéal engendré dans $A$ par $I$ ($\subset S\subset A$) et par $M^\ast$ ($\subset S$), [donc $J=\varepsilon^{-1}(I)$].
\end{minipage}

a) On peut trouver sur \deleted{$A$} $J\subset A$ une structure à p.d. telle que 1°) $S\longrightarrow A$ soit un hom. d'anneaux compatible aux p.d. sur $I\subset S$ et $J\subset A$ \quad 2°) Pour $(S,M^\ast)$ variable, cette structure à p.d. est fonctorielle. Ces conditions caractérisent de façon unique la structure à p.d. envisagée.

b) Soit $(B^\ast\deleted{\unsure{,K^\ast}})$ une $S$-algèbre graduée mod $2$, \deleted{alternée strict} alternée, munie d'un idéal $K^\ast$ gradué à p.d., \deleted{l'application de restriction} \inserted{compatibles avec les p.d.} sur $I\subset S$.
L'application de restriction
\[
\begin{aligned}
&\operatorname{Hom}_{S\text{-alg à p.d.}}\bigl((A(M^\ast),J(M^\ast)),(B^\ast,K^\ast)\bigr)\\
&\hspace{8em}\longrightarrow\operatorname{Hom}_{S\text{-mod grd}}(M^\ast,K^\ast)
\end{aligned}
\]
est bijective. (Donc il y a lieu d'appeler $A(M^\ast)$
la \unsure{$B$|$S$}-algèbre \inserted{graduée alternée} à p.d. \deleted{sur} $S$ \uline{engendrée} par le $S$-module gradué $M^\ast$.)

\uline{Dém} :
% archive sheet number: 214

\scanpage{216}
\begin{center}\uuline{POINTS DE GORENSTEIN D'UN MORPHISME}

par Manuela STOIA\end{center}

R\deleted{\unsure{e|a}}appelons qu'un anneau noethérien est de Gorenstein s'il est de dimension injective finie ( \handwritten{3} ). On va utiliser
le critère suivant: ``\uline{Un anneau local noetherien $\underline A$,quotient d'un anneau local noethérien régulier $\underline B$, est de Gorenstein
si et seulement si $\operatorname{Ext}^{\underline q}_{\underline B}(\underline A,\underline B)=0$ pour tous les entiers $\underline q$ à
l'exception d'un $\underline q_{\underline{\unsure{o|0}}}$ et $\operatorname{Ext}^{\handwritten{\underline q_{\underline o}}}_{\underline B}(\underline A,\underline B)\simeq\underline A$}''\handwritten{[(3),prop\unsure{.|:}1.3,ch 1,§ 1.4]}.

Soit $\underline k$ un corps algébriquement clos.On va considérer des variétés algébriques $\underline X$ sur $\underline k$ au sens de ( \handwritten{4} ), mais pouvant avoir des éléments nilpotents.
On dit que $\underline X$ est\unsure{:|.} de Gorenstein
au point $\underline x$ (réspect.de Gorenstein) si l'anneau
local $\handwritten{\mathscr O_{\underline X,\underline x}}$ est de Gorenstein (réspect.si tous les anneaux $\handwritten{\mathscr O_{\underline X,\underline x}}$ sont de Gorenstein).

\uline{Théorème}: \uline{Si $\underline f:\underline X\handwritten{\longrightarrow}\underline Y$ est un morphism\handwritten{e} plat de variétés algébriques, alors l'ensemble} :
\[
\underline{\underline G_{\handwritten{\underline f}}(\handwritten{\underline{\unsure{X|x}}})=\handwritten{\bigl\{}\handwritten{\underline x\in}\underline X/\underline X_{\handwritten{\underline f(\underline x)}}\quad\text{soit de Gorenstein au point }\handwritten{\underline x\bigr\}}}
\]
\uline{est ouvert dans $\underline X$ (pour $\underline y\handwritten{\in}\underline Y$,on note comme d'habitude par $\underline X_{\underline y}$
la variété algébrique- fibre d\deleted{\unsure{a}}e $\underline f$\deleted{\unsure{,}} au-dess\deleted{o}us de $\underline y$)}.

\uline{Démonstration}. L'énoncé est local sur $\underline X$,donc on peut
supposer $\underline X$ affine et soit $\deleted{\unsure{i|f}}\handwritten{\underline j}:\underline X\handwritten{\lhook\joinrel\longrightarrow}\handwritten{\underline k}^{\underline n}$ une immersion fermée.
Notons par $\underline i$ l'immersion donné par la composition
\diagram{\begin{tikzcd}[ampersand replacement=\&,column sep=large]
\underline X\arrow[r,"{\handwritten{\text{\textcircled{$\Gamma$}}_{\underline f}}}"]\&\underline X\times\underline Y\arrow[r,"{\handwritten{\underline j\times\underline 1_{\underline Y}}}"]\&\handwritten{\underline k}^{\underline n}\times\underline Y
\end{tikzcd}}
et par $\handwritten{\text{\textcircled{$\pi$}}}:\handwritten{\underline k}^{\underline n}\times\underline Y\handwritten{\longrightarrow}\underline Y$ la projection.On a la rélation $\underline f=\handwritten{\text{\textcircled{$\pi$}}}\circ\underline i$. Pour la commodité,
% archive sheet number: 215

\scanpage{217}
\[
A=A^0+A^1\supset J^0+J^1
\]
$A^0$ anneau comm

$A^1$ \quad $A^0$-module

$A^1\times A^1\xrightarrow{\varphi}A^0$ forme alternée (s'écrit $xy=\varphi(x,y)$)
\margin{$\bar\varphi:A^1\longrightarrow\operatorname{Hom}(A^1,A^0)$}

$J^0$ ideal de $A^0$

$J^1$ \deleted{$\unsure{J^0\cdot A^1}$} sous-module de $A^1$ tel que
\[
J^0 A^1\subset J^1\subset\bar\varphi^{-1}\bigl(\underset{\substack{\cap\\\operatorname{Hom}(A^1,A^0))}}{\underbrace{\operatorname{Hom}(A^1,J^0)}}\bigr)
\]
On suppose donnée une structure de p.d. sur $J^0\subset A^0$
\[
(A)\left\{\begin{aligned}
(0)\ &x^{(i)}\quad(i\geq0)\quad x^0=1,\quad x^{(1)}=x,\quad x^{(i)}\in J^0\quad(i\geq1)\\
(1)\ &(x+y)^{(i)}=\sum_{j+k=i}x^{(j)}y^{(k)}\qquad x,y\in J^0\qquad i\geq0\\
(2)\ &(\lambda^{\deleted{\unsure{(i)}}}x)^{(i)}=\lambda^i x^{(i)}\qquad\lambda\in A^0,\ x\in J^0\\
(3)\ &x^{(i)}x^{(j)}=c_{ij}x^{(i+j)}\qquad x\in J^0,\quad\deleted{\unsure{i}}i,j\geq0\quad c_{ij}=\frac{(i+j)!}{i!j!}\\
(4)\ &(x^{(i)})^{(j)}=d_{ij}x^{(ij)}\qquad x\in J^0\qquad i,j\geq\unsure{0|1}\quad d_{ij}=\frac{(ij)!}{(i!)^j j!}
\end{aligned}\right.
\]
\inserted{On suppose
\[
(\lambda_1^{\deleted{\unsure{+}}}x_1)^{(i)}=0\qquad i\geq2\qquad(\lambda_1\in A^1,x_1\in J^1).\tag*{(5)}
\]}
On prolonge les op $x^{(i)}$ de $J^0$ à \inserted{$J=J^0+$} $J^1$ par
\[
(x+u)^{(i)}=x^{(i)}+x^{(i-1)}u\qquad x\in J^0,\ u\in J^1,\ i\geq0\tag*{($\ast$)}
\]
\deleted{Voyons dans quelle mesure}

donc $\unsure{u|x}^{(i)}=0$ si $u\in J^1$, $i\geq2$

Voyons dans quelle mesure (A) \unsure{reste valable|restent valables} pour des éléments généraux de $J=J^0+J^1$

(0) OK (en convenant $x^{(i)}=0$ si $i<0$, $x\in J^0$) \inserted{car $J^0 A^1\subset J^1$}

(1) $\displaystyle(x+y)^{(i)}=\sum_{j+k=i}x^{(j)}y^{(k)}$ si $x,y\in J$ \uline{commutent} \uuline{strictement}
\deleted{(i.e. \unsure{posa|posant}} (i.e posant $x=x^0+x^1$, $y=y^0+y^1$, on veut (non seulmt $x^1y^1=y^1x^1$ \deleted{donc} ($=-x^1y^1$) i.e. $2x^1y^1=0$, mais même $x^1y^1=0$)

(2) $(\lambda x)^{(i)}=\lambda^i x^{(i)}$ si $\lambda\in A,x\in J$ comm\uline{utent} \uline{strictement} [\uline{grâce à} (5)]

(3) \deleted{OK} $x^{(i)}x^{(j)}=c_{ij}x^{(i+j)}$\qquad $\forall x\in J$\quad OK

(4) $(x^{(i)})^{(j)}=d_{ij}x^{(ij)}$
% archive sheet number: 216

\scanpage{218}\mspage{\handwritten{-4-}}
cation $\handwritten{\underline P'^{\underline q_{\underline{\unsure{o|0}}}}/\text{\textcircled{$\mathfrak m$}}_{\underline A}\underline P'^{\underline q_{\underline o}}\longrightarrow\underline P^{\underline q_{\underline o}}/\text{\textcircled{$\mathfrak m$}}_{\underline A}\underline P^{\underline q_{\underline o}}}$ étant injective,il en résulte que le complexe
\[
\handwritten{\cdots\longrightarrow\underline P^{\underline q_{\underline o}-2}\longrightarrow\underline P^{\underline q_{\underline o}-1}\longrightarrow\underline P'^{\underline q_{\underline o}}\longrightarrow0\longrightarrow\cdots}\ .
\]
vérifie les conditions du lemme. On peut donc supposer dans l'énoncé que $\handwritten{\underline P^{\underline q}=0}$ pour $\underline q>\underline q_{\underline o}$.Utilisons [(\handwritten{1}), Ch.III,§.5] pour montrer que $\underline H$ est $\underline A$-plat.Le complexe $\underline P^\bullet$ est une résolution pour $\underline H$ par des $\underline A$-modules plats.Par consequence,les modules
$\handwritten{\operatorname{Tor}^{\underline A}_\bullet(\underline A/\text{\textcircled{$\mathfrak m$}}_{\underline A},\underline H)}$ sont les modules de cohomologie du complexe $\handwritten{\underline P^{\unsure{\circ|\bullet}}\otimes_{\underline A}\underline A/\text{\textcircled{$\mathfrak m$}}_{\underline A}=\underline P^\circ/\text{\textcircled{$\mathfrak m$}}_{\underline A}\underline P^\circ}$ et la conclusion résulte.

\uline{Corollaire}. \uline{Soit $\underline f:\underline X\handwritten{\longrightarrow}\underline Y$ un morphisme propre et plat de variétés algébriques. Alors l'ensemble}
\[
\underline{\handwritten{\underline G_{\underline f}(\underline Y)=\{\underline y\in\underline Y\mid\underline X_{\underline y}}\quad\text{soit de Gorenstein}\ \handwritten{\}}}
\]
\uline{est ouvert dans $\underline Y$.}

\begin{flushright}L'Institut Central de Mathématique\\
14,\deleted{\unsure{x|r}}rue de l'Académie , Bucarest\\
\handwritten{ROUMANIE}\end{flushright}

\noindent\begin{minipage}[t]{.33\linewidth}
\diagram{\(\begin{array}{cl|}\text{---}&\handwritten{\text{latine}}\\
\handwritten{\text{\tikz[baseline=-.4ex]\draw[densely dotted] (0,0) circle (.63em);}}&\handwritten{\text{gothique}}\\
\handwritten{\bigcirc}&\handwritten{\text{grec}}\\
\handwritten{\text{\tikz[baseline=-.4ex]\draw[black!25] (0,0) circle (.63em);}}&\handwritten{\text{ronde}}\\ \hline
\end{array}\)}
\end{minipage}\hfill\begin{minipage}[t]{.63\linewidth}
\noindent\rule{\linewidth}{.3pt}\par
\handwritten{\uline{Bibliographie}}

\handwritten{[1] Bourbaki, N. -- Algèbre commutative, Paris, Hermann 1961}

\handwritten{[2] Grothendieck, A -- E.G.A., IV , 1965}

\handwritten{[3] Samuel, P -- Séminaire d'Algèbre commutative, 1966/67.}

\handwritten{[4] J. P. Serre -- Faisceaux algébriques cohérents, Ann. Math. 1955}
\end{minipage}
% archive sheet number: 217

\scanpage{219}
\uline{Quelques complexes d'\unsure{Algèbres|algèbres|algèbre} différentielles.}

$A$ anneau commutatif, $M$ $A$-module. \margin{On considère $\mathrm{Sym}^{\ast}(M)\otimes\bigwedge^{\ast}M$ et $\Gamma^{\ast}(M)\otimes\bigwedge^{\ast}M$ \uline{comme} $A$-\uline{alg} \uline{graduées par degré de} $\bigwedge^{\ast}M$}

\margin{De Rham}% slanted
\textcircled{1} Sur $\mathrm{Sym}^{\ast}(M)\otimes\bigwedge^{\ast}M$ \deleted{(algèbre bigraduée)}, il y a un
unique opérateur $d$ qui soit de \deleted{a) $A$-linéaire} \inserted{$A$-algèbre gr. anticommutative, une} \uline{dérivation} \deleted{(pour $\mathrm{Sym}^{\ast}M$ de degré pair et \deleted{$\mathrm{Sym}$} $\bigwedge^{\ast}M$ de degré impair)}
b) qui soit nulle sur $\bigwedge^{1}M$, et qui
sur $\mathrm{Sym}^{1}(M)$ soit l'isom. can. $\mathrm{Sym}^{1}(M)\longrightarrow\bigwedge^{1}M$.

\deleted{\unsure{Il|Le}} Ce $d$ \deleted{\unsure{est}} satisfait $d^2=0$, et applique $\mathrm{Sym}^{p}\otimes\bigwedge^{q}$ dans $\mathrm{Sym}^{p-1}\otimes\bigwedge^{q+1}$. \deleted{Ce $d$ est} \inserted{Il est $\bigwedge^\ast M$-linéaire.}
C'est aussi le complexe de De Rham de $\mathrm{Sym}^{\ast}(M)$ sur $A$, donc le complexe de De Rham du fibré vectoriel $V(M)$ sur $A$.
\deleted{Si \deleted{\unsure{tf}} $M$ est libre et
Si $M$ est libre,}

Supposons $M$ plat\margin{inutile}, \deleted{(par ex., plat tf)} $A$ de car. 0 (i.e. $A$ une $\mathbb Q$-algèbre), alors
\deleted{$H^{\ast}(\mathrm{Sym}^{\ast}(M)\otimes\bigwedge^{\ast}M)$ (gradué par le degré de $\bigwedge^{\ast}M$) est nul en degrés $>0$, et en degré} (graduant suivant le degré de $\bigwedge^\ast M$
\[
H^i_{\mathrm{DR}}(\mathrm{Sym}^{\ast}(M)\otimes\bigwedge^{\ast}(M))\simeq
\begin{cases}0&\text{si }i\ne0\\
\underset{\substack{\shortparallel\\[-2pt]\mathrm{Sym}^{0}\otimes\bigwedge^{0}}}{A}&\text{si }i=0\end{cases}
\]
Si $A$ est de car. $p>0$\inserted{(i.e. est une $\mathbb F_p$-alg), $M$ plat}, alors on a \unsure{.....|....|\ldots} (théorie des opérations de Cartier ----)

\margin{Koszul}% slanted
\textcircled{2} Sur $\mathrm{Sym}^{\ast}(M)\otimes\bigwedge^\ast M$ \deleted{(algèbre bigraduée)} il y a un unique opérateur $\delta$ qui est \deleted{a) $A$-linéaire}
a) une \inserted{anti}dérivation \deleted{si} de \inserted{$A$-alg. graduées, et} \deleted{on prend $\mathrm{Sym}^{1}M$ de degré pair, et b) $\bigwedge^{\ast}M$ de degré impair,}
b) nulle sur \deleted{$\bigwedge^{\unsure{q|0|\ast}}M$} $\mathrm{Sym}^{1}M$, et sur $\bigwedge^{1}M$ c'est l'isom can. $\bigwedge^{1}M\xrightarrow{\sim}\mathrm{Sym}^{1}M$.

Alors $\delta$ satisfait $\delta^2=0$, il applique $\mathrm{Sym}^{p}\otimes\bigwedge^{q}$ dans $\mathrm{Sym}^{p+1}\otimes\bigwedge^{q-1}$. Si $M$ est libre (plus gén, plat) et sans condition de car. sur $A$, on a (\deleted{gr} graduant la coh. par le degré de $\bigwedge^{\ast}$ \inserted{qu'on écrit en bas}
(de sorte que $\delta$ est de degré \deleted{$-1$}\inserted{inf. $1$})):
\[
H^{\mathrm{Kos}}_{\unsure{q|i}}(\mathrm{Sym}^{\ast}M\otimes\bigwedge^\ast M)=
\begin{cases}0&\text{si }i\ne0\\
\simeq\underset{\substack{\shortparallel\\[-2pt]\mathrm{Sym}^{0}\otimes\bigwedge^{0}}}{A}&\text{'' }i=0\end{cases}.
\]
% archive sheet number: 218
